\documentclass[a4paper,twoside]{article}

\usepackage{morewrites}

\usepackage[svgnames,dvipsnames,x11names]{xcolor}
\usepackage{graphicx}
\usepackage[english]{babel}
\usepackage[colorlinks=true, linkcolor=NavyBlue, urlcolor=NavyBlue, citecolor=NavyBlue]{hyperref}
\usepackage[a4paper]{geometry}
\usepackage{ts-fonts}
\usepackage{tcolorbox}
\usepackage{fvextra}
\usepackage{amsmath}
\usepackage{float}
\usepackage{seqsplit}
\usepackage{ts-code}

\usepackage{lastpage}
\usepackage{refcount}
\makeatletter
\def\ps@plain{
\let\@oddhead\@empty
\let\@evenhead\@empty
\def\@oddfoot{
\hfil
\ifnum\getpagerefnumber{LastPage}>1\relax
\thepage
\fi
\hfil}
\let\@evenfoot\@oddfoot
}
\makeatother

\title{TeXSmith}
\author{}\date{}
\begin{document}
\maketitle
And of course, the grand finale---the true climax of the project---is that this very documentation can itself be converted into a \LaTeX{} document using TeXSmith.
\begin{tscode}[lang=bash, engine=pygments]
\begin{Verbatim}[commandchars=\\\{\},breaklines, breakanywhere, commandchars=\\\{\}]
git\PY{+w}{ }clone\PY{+w}{ }https://github.com/yves\PYZhy{}chevallier/texsmith.git
\PY{n+nb}{cd}\PY{+w}{ }texsmith
uv\PY{+w}{ }sync\PY{+w}{ }\PYZhy{}\PYZhy{}with\PY{+w}{ }docs
\PY{n+nb}{export}\PY{+w}{ }\PY{n+nv}{TEXSMITH\PYZus{}MKDOCS\PYZus{}BUILD}\PY{o}{=}\PY{l+m}{1}\PY{+w}{ }\PY{c+c1}{\PYZsh{} Enable PDF build}
uv\PY{+w}{ }run\PY{+w}{ }texsmith\PY{+w}{ }mkdocs\PY{+w}{ }build
\end{Verbatim}
\end{tscode}
You can download it directly from the \href{https://github.com/yves-chevallier/texsmith/releases}{release page}.
\end{document}
